\subsection{广义 Calderón--Zygmund 算子}
\label{sec:ch5-generalized-calderon-zygmund}

到目前为止, 所研究的算子都可写成与缓增分布的卷积. 这一假设的实际优点是可以用 Fourier 变换推出 $L^2$ 有界性. 一旦假设 $L^2$ 有界, 推广到其余 $L^p$ 空间只依赖于 Hörmander 条件, 而该条件可以调整到非卷积算子.

\hypertarget{chapter-05-page-111}{}
令 $\Delta=\{(x,x):x\in\mathbb R^n\}$ 是 $\mathbb R^n\times\mathbb R^n$ 的对角线. 用几乎与定理~\ref{thm:ch5-calderon-zygmund} 相同的证明可得下述结论.

\begin{Theorem}
\label{thm:ch5-generalized-calderon-zygmund-bounds}
设 $T$ 是 $L^2(\mathbb R^n)$ 上的有界算子, $K$ 是 $\mathbb R^n\times\mathbb R^n\setminus\Delta$ 上的函数, 并且对每个具有紧支集的 $f\in L^2(\mathbb R^n)$,
\[
 Tf(x)=\int_{\mathbb R^n}K(x,y)f(y)\,dy,
 \qquad x\notin\operatorname{supp}f.
\]
再假设 $K$ 满足
\begin{equation}
\label{eq:ch5-generalized-hormander-y}
 \int_{|x-y|>2|y-z|}|K(x,y)-K(x,z)|\,dx\le C,
\end{equation}
\begin{equation}
\label{eq:ch5-generalized-hormander-x}
 \int_{|x-y|>2|x-w|}|K(x,y)-K(w,y)|\,dy\le C.
\end{equation}
则 $T$ 是弱型 $(1,1)$ 的, 并且是强型 $(p,p)$ 的, $1<p<\infty$.
\end{Theorem}

上述核表示只对具有紧支集的 $L^2$ 函数提出, 但证明中它只用于 Calderón--Zygmund 分解所得的坏函数 $b_j$, 因而已经足够. 条件 \eqref{eq:ch5-generalized-hormander-y} 用于证明 $T$ 的弱型 $(1,1)$ 不等式, 条件 \eqref{eq:ch5-generalized-hormander-x} 用于证明伴随算子的相应不等式.

沿用 Coifman 与 Meyer \MBUnresolvedReference{citation}{[2]} 的术语, 若存在 $\delta>0$ 使
\begin{equation}
\label{eq:ch5-standard-kernel-size}
 |K(x,y)|\le\frac C{|x-y|^n},
\end{equation}
\begin{equation}
\label{eq:ch5-standard-kernel-y-regularity}
 |K(x,y)-K(x,z)|\le C\frac{|y-z|^\delta}{|x-y|^{n+\delta}},
 \qquad |x-y|>2|y-z|,
\end{equation}
\begin{equation}
\label{eq:ch5-standard-kernel-x-regularity}
 |K(x,y)-K(w,y)|\le C\frac{|x-w|^\delta}{|x-y|^{n+\delta}},
 \qquad |x-y|>2|x-w|,
\end{equation}
则称 $K:\mathbb R^n\times\mathbb R^n\setminus\Delta\to\mathbb C$ 为标准核. 标准核显然满足 \eqref{eq:ch5-generalized-hormander-y} 与 \eqref{eq:ch5-generalized-hormander-x}.

下面给出具有这类核的算子例子.

第一类是 Lipschitz 曲线上的 Cauchy 积分. 设 $A$ 是 $\mathbb R$ 上的 Lipschitz 函数, $A'=a\in L^\infty$, 并令 $\Gamma=(t,A(t))$. 对 $f\in\mathcal S(\mathbb R)$,
\[
 C_\Gamma f(z)=\frac1{2\pi i}\int_{-\infty}^\infty
 \frac{f(t)(1+ia(t))}{t+iA(t)-z}\,dt
\]
在 $\Omega_+=\{z=x+iy:y>A(x)\}$ 上定义解析函数. 它在 $\Gamma$ 上的边界值导出算子
\[
 Tf(x)=\lim_{\varepsilon\to0}\int_{|x-y|>\varepsilon}
 \frac{f(y)}{x-y+i(A(x)-A(y))}\,dy,
\]
其核
\begin{equation}
\label{eq:ch5-cauchy-lipschitz-kernel}
 K(x,y)=\frac1{x-y+i(A(x)-A(y))}
\end{equation}
以 $\delta=1$ 满足 \eqref{eq:ch5-standard-kernel-y-regularity} 与 \eqref{eq:ch5-standard-kernel-x-regularity}.

\hypertarget{chapter-05-page-112}{}
第二类是 Calderón 交换子. 当 $\|A'\|_\infty<1$ 时, 可将 \eqref{eq:ch5-cauchy-lipschitz-kernel} 展为几何级数. 因而对 Lipschitz 函数 $A$ 和整数 $k\ge0$, 自然要考察
\[
 T_kf(x)=\lim_{\varepsilon\to0}\int_{|x-y|>\varepsilon}
 \left(\frac{A(x)-A(y)}{x-y}\right)^k\frac{f(y)}{x-y}\,dy.
\]
相应的核 $K_k(x,y)=((A(x)-A(y))/(x-y))^k(x-y)^{-1}$ 也是 $\delta=1$ 的标准核.

\begin{Definition}
\label{def:ch5-calderon-zygmund-operator}
若算子 $T$ 满足:
\begin{enumerate}
\item $T$ 在 $L^2(\mathbb R^n)$ 上有界;
\item 存在标准核 $K$, 使对每个具有紧支集的 $f\in L^2$,
\[
 Tf(x)=\int_{\mathbb R^n}K(x,y)f(y)\,dy,
 \qquad x\notin\operatorname{supp}f,
\]
\end{enumerate}
则称 $T$ 为一个广义 Calderón--Zygmund 算子.
\end{Definition}

定理~\ref{thm:ch5-generalized-calderon-zygmund-bounds} 表明 Calderón--Zygmund 算子在 $L^p$, $1<p<\infty$, 上有界并且是弱型 $(1,1)$ 的. 因而, 给定具有标准核的算子后, 问题归结为证明其 $L^2$ 有界性. 我们将在\MBUnresolvedReference{chapter}{第 9 章}回到这个问题.
